\section{이전 자료의 교정과 단순화}\label{LEAD:corrections}
이 절은 Round~4에서 이전 두 문서의 진술을 다시 검토하여 얻은 교정과 단순화를 모은다. 각 항목의 근거는 괄호 안의 절이다.

\begin{enumerate}[leftmargin=1.6em, itemsep=4pt]
\item \textbf{cascade 조립정리의 서로소 조건은 $n=\Lam_L$에서 불필요하다.} \PROVED{}
P1--P4 보고서 정리 8.10은 단계 족 $D_0,\dots,D_K$가 정수 집합으로서 쌍마다 서로소일 것을 요구하고, 반례 8.12($n=6$, $4=2+2$)로 이 조건이 필요하다고 기록했다. 그러나 $n=\Lam_L$이면 각 단계의 출력들을 다중집합으로 합친 뒤 중복 제거 정리(STRUCT 절)를 적용하면 항수를 늘리지 않고 서로 다른 진약수의 표현을 얻는다. 합이 $a<\Lam_L$이므로 정리의 가정이 충족된다. 반례 8.12 자체도 $\Lam_3=6$에서 $d=2$의 최소 부족 소수가 $p=3$, $r=2$이므로 $2+2=3+1$로 바뀌어 $4/6=1/2+1/6$이 된다. 따라서 반례 8.12는 ``단계 출력을 그대로 이어 붙이는 조립''에 대한 반례일 뿐, 결론 $H(L)\le\sum_k s_k$에 대한 반례가 아니다. 일반 정수 $n$에서는 이 교정이 성립하지 않는다($n=10$, $4=2+2$는 $\{1,2,5\}$의 서로 다른 부분합이 아니다).
\item \textbf{라벨 분리 설계(P1--P4 명제 8.21)와 Newton 보정(Baker 12.3절)은 불필요하다.} \PROVED{} 같은 이유로, 단계 사이의 약수 충돌이나 네 항 표현의 중복 선택을 따로 제거할 필요가 없다. 반복을 허용한 순서쌍의 개수만 세면 된다.
\item \textbf{P1b 주의 5.19(c)의 발견적 경로.} P1--P4 보고서는 Vinogradov--Korobov 비공명과 국소극한정리로 $H(L)\ll L/(\log L)^{3/2-\eps}$를 얻을 수 있을 것이라고 \HEUR{}로 기록했다. 본 라운드의 VKGAP 절이 이 경로를 증명으로 바꾸는 작업을 수행했다(판정은 해당 절과 감사 절). 리드의 검토에서 한 가지 함정이 발견되었다: 풀 $P=(L/2,L]$에서 포함 여부를 독립 Bernoulli로 고르면 $S=\sum\log p$가 간격 $\log L$의 격자에 가깝게 모여, 주파수 $2\pi/\log L$ 부근에서 특성함수가 $\exp(-cn/(\log L)^2)$ 이상으로 남는다. 크기를 고정한 $n$-부분집합(또는 $|E|=n$으로 조건화한 2차원 국소극한)을 쓰면 $S=n\log L-W$가 되어 이 문제가 사라진다.
\item \textbf{다음 연구 목표 \#1(합성 법 혼합, CMM).} P1--P4 보고서 12.4절이 \OPEN{}으로 남긴 CMM은 CMM 절에서 다루었다(범위와 판정은 해당 절).
\end{enumerate}
